This thesis addresses the automatic construction of knowledge graphs from Slovak legal documents using large language models. Classical vector-based RAG fails on legislation for three reasons: paragraphs cross-reference one another throughout the document, and tables and formulas fall apart during plain text conversion.

The thesis designs and implements a pipeline that solves all three problems. Preprocessing extracts tables and formulas through separate routes. A three-step extraction first collects entity and relation types from the document in an open-domain setting, then unifies them into a coherent schema using a larger model, and finally uses this schema to drive the extraction of concrete instances. The resulting nodes and edges are loaded into a Neo4j graph database with three vector indexes over chunks, entities and relations to support hybrid retrieval.

The solution is evaluated on Act No. 222/2004 Coll. on Value Added Tax. Four combinations of model and prompting regime are compared against a manually annotated reference using F1 score, hallucination rate, omission rate and normalized graph edit distance. A second set of experiments compares answers from classical vector RAG with those produced by GraphRAG over the same act, and assesses whether the proposed approach more accurately retrieves the supporting provisions.